\begin{table}[t]
\caption{Multi-election robustness. Fraction of statewide races under which
each vote-based statistic flags the enacted plan as an outlier (beyond the
99th ensemble percentile), against the topology-only detector, whose verdict
consumes no vote data and so does not vary with the election.}
\label{tab:elections}
\centering
\footnotesize
\begin{tabular}{lcccc c}
\toprule
 & mean & efficiency & & & geometry \\
State & median & gap & declin. & lean & $p$ \\
\midrule
NC (R) & 5/5 & 4/5 & 4/5 & 0.49 & 0.36 \\
PA (R) & 4/9 & 3/9 & 4/9 & 0.52 & \textbf{0.009} \\
MD (D) & 0/6 & 2/6 & 4/6 & 0.55 & \textbf{0.016} \\
\bottomrule
\end{tabular}

\vspace{2pt}
{\scriptsize Fractions are flagged races over available races. ``lean'' is the
mean statewide two-party Democratic share. ``geometry $p$'' is the topology-only
detector's $\varepsilon$-outlier $p$-value (one value per state, election-invariant);
bold marks $p<0.05$. Mean-median flags the clear Republican gerrymander (NC)
everywhere and the lopsided Democratic gerrymander (MD) nowhere.}
\end{table}
